% TEMPLATE-MILC-v3.tex - plantilla maestra para todas las guias MILC v3
% Ver scripts/generadores/build-guia-milc.py para la lista completa de
% claves esperadas. El script valida que no queden placeholders sin
% reemplazar antes de escribir el archivo final.

\documentclass[11pt,letterpaper]{article}
\usepackage[margin=1.75cm]{geometry}
\usepackage[table]{xcolor}
\usepackage{fontspec}
\usepackage[spanish]{babel}
\usepackage{tikz}
\usepackage[most]{tcolorbox}
\usepackage{tabularx}
\usepackage{array}
\usepackage{fancyhdr}
\usepackage{titlesec}
\usepackage{ragged2e}
\usepackage{enumitem}
\usepackage{microtype}

\setmainfont{Helvetica Neue}
\setsansfont{Helvetica Neue}
\setlength{\parindent}{0pt}
\setlength{\parskip}{6pt}
\hyphenpenalty=200
\exhyphenpenalty=200
\pretolerance=400
\tolerance=2000
\setlength{\emergencystretch}{1em}
\renewcommand{\arraystretch}{1.25}
\setlist{leftmargin=5mm,itemsep=1mm,topsep=1mm}
\newcolumntype{Y}{>{\RaggedRight\arraybackslash}X}

\definecolor{milcMagenta}{HTML}{E6007E}
\definecolor{milcVino}{HTML}{5A0038}
\definecolor{milcTurquesa}{HTML}{0093A5}
\definecolor{milcVerde}{HTML}{58A923}
\definecolor{milcMostaza}{HTML}{E5B400}
\definecolor{milcOcre}{HTML}{B86600}
\definecolor{milcNegro}{HTML}{241816}
\definecolor{milcGris}{HTML}{F7F7F4}
\definecolor{milcLinea}{HTML}{E8E2DD}
\definecolor{milcGradePrimary}{HTML}{0066FF}
\definecolor{milcGradeDark}{HTML}{003D99}
\definecolor{milcGradeSoft}{HTML}{D6E8FF}
\definecolor{milcGradeAccent}{HTML}{CFE7FF}
\definecolor{milcGradeText}{HTML}{FFFFFF}
\color{milcNegro}

\pagestyle{fancy}
\fancyhf{}
\lhead{\small Tecnología e Informática · Grado 6}
\rhead{\small Guía 14 · MILC v3}
\cfoot{\small\thepage}
\renewcommand{\headrulewidth}{0pt}

\newtcolorbox{titlebox}[1]{
  enhanced,colback=#1,colframe=#1,arc=7mm,boxrule=0pt,
  left=7mm,right=7mm,top=3mm,bottom=3mm,
  before skip=8pt,after skip=8pt,
  fontupper=\sffamily\bfseries\Large\color{white}
}

\newtcolorbox{softbox}[3]{
  enhanced,colback=#2,colframe=#1,borderline west={4pt}{0pt}{#1},
  arc=2mm,boxrule=.4pt,left=4mm,right=4mm,top=3mm,bottom=3mm,
  before skip=6pt,after skip=8pt,title={#3},coltitle=white,
  colbacktitle=#1,fonttitle=\sffamily\bfseries,fontupper=\RaggedRight,
  attach boxed title to top left={xshift=3mm,yshift=-2mm},
  boxed title style={arc=2mm, boxrule=0pt}
}

\newtcolorbox{drawbox}[2]{
  enhanced,colback=white,colframe=#1,arc=4mm,boxrule=1pt,
  left=5mm,right=5mm,top=4mm,bottom=4mm,before skip=8pt,after skip=8pt,
  title={#2},coltitle=white,colbacktitle=#1,fonttitle=\sffamily\bfseries,
  fontupper=\RaggedRight,
  attach boxed title to top left={xshift=4mm,yshift=-2mm},
  boxed title style={arc=3mm, boxrule=0pt}
}

\newtcolorbox{infoband}[1]{
  enhanced,colback=#1!8,colframe=#1!50,arc=2mm,boxrule=.5pt,
  left=4mm,right=4mm,top=2mm,bottom=2mm,before skip=4pt,after skip=4pt,
  fontupper=\small\RaggedRight
}

% Puente narrativo entre secciones (contrato editorial Conéctate).
% Da continuidad: "Acabas de X. Ahora vas a Y porque Z."
% Visualmente: banda izquierda en color del grado, texto en itálica pequeña.
\newtcolorbox{puentebox}{
  enhanced,colback=milcGradeSoft,colframe=milcGradeDark,
  borderline west={3pt}{0pt}{milcGradeDark},
  arc=0mm,boxrule=0pt,left=5mm,right=4mm,top=2.5mm,bottom=2.5mm,
  before skip=4pt,after skip=8pt,
  fontupper=\itshape\small\color{milcNegro}\RaggedRight
}
\newcommand{\puente}[1]{%
  \begin{puentebox}%
    \textcolor{milcGradeDark}{\textbf{\textrightarrow}}\hspace{2mm}#1%
  \end{puentebox}%
}

\newcommand{\linead}[1]{\noindent\color{milcLinea}\rule{#1}{0.5pt}\color{milcNegro}}
\newcommand{\respuesta}[1]{\par\vspace{2mm}\foreach \n in {1,...,#1}{\noindent\color{milcLinea}\rule{\linewidth}{0.5pt}\par\vspace{5mm}}\color{milcNegro}}
\newcommand{\checkbox}{\textcolor{milcGradeDark}{\fbox{\phantom{x}}}\hspace{5pt}}

\newcommand{\phasecard}[4]{
\begin{tcolorbox}[enhanced,colback=#3,colframe=#2,arc=3mm,boxrule=.5pt,height=3.05cm,valign=center,halign=center,left=1.5mm,right=1.5mm,top=2mm,bottom=2mm]
{\sffamily\bfseries\fontsize{22}{24}\selectfont\textcolor{#2}{#1}}\par
{\sffamily\bfseries\small #4}
\end{tcolorbox}
}

\begin{document}
\thispagestyle{empty}
\begin{tikzpicture}[remember picture,overlay]
  \fill[white] (current page.south west) rectangle (current page.north east);
  \path[fill=milcGradePrimary,rounded corners=16mm]
    ([xshift=.55cm,yshift=-.55cm]current page.north west) rectangle
    ([xshift=-.55cm,yshift=3.65cm]current page.south east);

  \node[anchor=north west,text=milcGradeText] at ([xshift=2.0cm,yshift=-1.85cm]current page.north west)
    {\sffamily\bfseries\fontsize{14}{16}\selectfont GRADO};
  \node[anchor=north west,text=milcGradeText,opacity=.82] at ([xshift=2.0cm,yshift=-2.75cm]current page.north west)
    {\sffamily\bfseries\fontsize{9}{11}\selectfont SERIE GUÍAS · TECNOLOGÍA E INFORMÁTICA};

  \begin{scope}[shift={([xshift=-3.05cm,yshift=-2.55cm]current page.north east)}]
    \fill[white,opacity=.80] (-.62,-.52) rectangle (.62,.52);
    \draw[milcGradeText,line width=1pt] (-.62,-.52) rectangle (.62,.52);
    \fill[milcGradeDark] (-.42,-.38) rectangle (-.22,.06);
    \fill[milcGradeAccent] (-.10,-.38) rectangle (.10,.24);
    \fill[milcGradePrimary!60!black] (.22,-.38) rectangle (.42,.38);
  \end{scope}

  \node[anchor=west,text=milcGradeText] at ([xshift=1.9cm,yshift=-12.85cm]current page.north west)
    {\sffamily\bfseries\fontsize{124}{124}\selectfont 6};
  \draw[milcGradeText,line width=1.7pt]
    ([xshift=4.52cm,yshift=-11.68cm]current page.north west) circle (.13cm);

  \node[anchor=west,text=milcGradeText,text width=8.7cm,align=left] at ([xshift=10.35cm,yshift=-11.6cm]current page.north west)
    {
     {\sffamily\bfseries\fontsize{19}{22}\selectfont Periféricos ---\\las manos,\\los ojos y\\la voz del\\computador}\\[4mm]
     {\sffamily\bfseries\fontsize{23}{26}\selectfont Guía 14-6-TIC}\\[2mm]
     {\sffamily\fontsize{13}{16}\selectfont Período 2 · Hardware y software}\\[5mm]
     {\sffamily\bfseries\fontsize{11}{13}\selectfont Institución Educativa Sor María Juliana}\\[1mm]
     {\sffamily\fontsize{10}{12}\selectfont Autor: Dr. Álvaro Cárdenas Orozco}
    };

  \path[fill=milcGradeSoft,rounded corners=7mm]
    ([xshift=1.35cm,yshift=.85cm]current page.south west) rectangle
    ([xshift=-1.35cm,yshift=4.25cm]current page.south east);
  \draw[milcGradeDark,line width=.65pt,rounded corners=7mm]
    ([xshift=1.35cm,yshift=.85cm]current page.south west) rectangle
    ([xshift=-1.35cm,yshift=4.25cm]current page.south east);
  \node[anchor=center,text=milcNegro,text width=17.0cm,align=center] at ([yshift=2.55cm]current page.south)
    {\sffamily\fontsize{10}{12}\selectfont Metodología MILC · Apertura ancestral · Pensamiento computacional · Triángulo Dussel-Estoicismo-Floridi\\
    \textcolor{milcGradeDark}{\bfseries Producto:} Una \textbf{tabla maestra de 12 periféricos} en una hoja del cuaderno, con 4 columnas: \emph{nombre}, \emph{tipo} (entrada / salida / mixto), \emph{cómo se conecta} (USB, cable, Bluetooth, inalámbrico), \emph{1 ejemplo de uso real}. Más \textbf{un diagrama lateral} donde dibujas un computador en el centro y \textbf{6 periféricos conectados a él con flechas} (3 de entrada apuntando al computador, 3 de salida apuntando hacia afuera). Firma de cierre: \emph{``Conozco las manos, los ojos y la voz de mi computador.''}.};
\end{tikzpicture}
\mbox{}
\newpage

\section*{Apertura: Periféricos --- las manos, los ojos y la voz del computador}

\begin{softbox}{milcGradeDark}{milcGradeSoft}{Saber ancestral}
\textbf{En el taller de don Lucho el relojero de Cartago, cada herramienta tenía su lugar y su uso.} Si entrabas a su vitrina de la calle 14, veías colgadas en orden las herramientas de su \textbf{oficio}: la \textbf{lupa} para ver lo pequeño, las \textbf{pinzas largas} para agarrar tornillos diminutos, los \textbf{destornilladores en miniatura} para abrir tapas, el \textbf{martillito de bronce} para ajustar suavemente, el \textbf{paño de gamuza} para limpiar sin rayar, el \textbf{aceite especial} para lubricar engranajes. \textbf{Cada herramienta servía para una cosa específica y don Lucho sabía cuándo sacar cuál}. La lupa no servía para apretar tornillos, ni el martillo para mirar dentro. Una persona inexperta usaba la herramienta equivocada y dañaba el reloj. \textbf{El relojero sabio sabía cuál herramienta usar cuándo}. Hoy el computador es la máquina: las herramientas que la rodean se llaman \textbf{periféricos}. Cada una tiene una función específica, y conocerlas te hace usuario sabio.
\end{softbox}

\begin{softbox}{milcTurquesa}{milcGris}{Saber contemporáneo}
Un \textbf{periférico} es cualquier pieza que se conecta \emph{por fuera} del computador. Vienen en 3 grandes tipos. \textbf{Entrada}: por donde tú le das información al computador (teclado, ratón, micrófono, cámara, escáner, joystick, lápiz óptico). \textbf{Salida}: por donde el computador te muestra resultados (pantalla, impresora, altavoz, audífonos, proyector, vibrador). \textbf{Mixtos (entrada Y salida a la vez)}: la pantalla táctil, los audífonos con micrófono, el USB, el módem, el monitor táctil. Hay también una clasificación por cómo se \textbf{conectan}: con cable (USB, HDMI, jack 3.5mm) o sin cable (Bluetooth, wifi, infrarrojo). Reconocer todo esto te permite \textbf{escoger el periférico correcto para lo que necesites} y resolver problemas básicos sin depender de un técnico.
\end{softbox}

\begin{softbox}{milcMostaza}{milcGris}{Pregunta-puente}
\textit{Cuando ves a alguien jugar Roblox en computador, usa al menos \textbf{4 periféricos al mismo tiempo}. ¿Sabrías nombrarlos? Pista: 2 son de entrada y 2 de salida.}
\end{softbox}

\begin{softbox}{milcVerde}{milcGris}{Saber y saber hacer}
Al terminar la clase: (1) podrás \textbf{identificar} 12 periféricos comunes; (2) sabrás \textbf{analizar} si cada uno es entrada, salida o mixto; (3) podrás \textbf{distinguir} las formas de conexión (USB, HDMI, jack, Bluetooth); (4) habrás \textbf{creado} tu tabla maestra de periféricos + diagrama.
\end{softbox}

\small
\begin{tabularx}{\textwidth}{>{\bfseries}p{3.0cm}Y}
\rowcolor{milcGradePrimary!12}Autor & Dr. Álvaro Cárdenas Orozco\\
DBA & Identifica componentes y sistemas tecnológicos en su entorno (MEN, T\&I 6°)\\
Referentes & Saber ancestral del relojero · Sistemas como cadena de oficios · Diagnóstico paso a paso\\
Producto final & Una \textbf{tabla maestra de 12 periféricos} en una hoja del cuaderno, con 4 columnas: \emph{nombre}, \emph{tipo} (entrada / salida / mixto), \emph{cómo se conecta} (USB, cable, Bluetooth, inalámbrico), \emph{1 ejemplo de uso real}. Más \textbf{un diagrama lateral} donde dibujas un computador en el centro y \textbf{6 periféricos conectados a él con flechas} (3 de entrada apuntando al computador, 3 de salida apuntando hacia afuera). Firma de cierre: \emph{``Conozco las manos, los ojos y la voz de mi computador.''}.\\
\end{tabularx}
\normalsize

\puente{Hoy haces 4 cosas: reconoces 12 periféricos, los clasificas en entrada/salida/mixtos, ves cómo se conectan, y armas tu tabla maestra.}

\begin{titlebox}{milcGradeDark}
Ruta MILC: así vamos a aprender
\end{titlebox}
\begin{tabularx}{\textwidth}{>{\centering\arraybackslash}X>{\centering\arraybackslash}X>{\centering\arraybackslash}X>{\centering\arraybackslash}X}
\phasecard{1}{milcVerde}{milcGris}{Escuta\\[-1mm]\small Observo, escucho y pregunto.} &
\phasecard{2}{milcTurquesa}{milcGris}{Sistematización\\[-1mm]\small Pensamiento computacional.} &
\phasecard{3}{milcMagenta}{milcGris}{Praxis\\[-1mm]\small Construyo y aplico.} &
\phasecard{4}{milcVino}{milcGris}{Evaluación\\[-1mm]\small Reflexión liberadora.}
\end{tabularx}


\puente{Empezamos identificando 12 periféricos en una lista mezclada.}

\begin{titlebox}{milcVerde}
Fase 1 · Escuta
\end{titlebox}

\begin{softbox}{milcVerde}{milcGris}{Escena para escuchar}
\textbf{Actividad 1 · IDENTIFICA --- Reconoce los 12 periféricos} (12 min · individual). En el cuaderno, escribe esta lista numerada del 1 al 12: \emph{(1) teclado, (2) monitor, (3) ratón (mouse), (4) impresora, (5) cámara web, (6) audífonos con micrófono, (7) altavoz, (8) micrófono solo, (9) joystick de juegos, (10) pantalla táctil, (11) escáner, (12) memoria USB}. Al lado de cada uno, en una sola línea, \textbf{escribe para qué crees que sirve} con tus palabras. Si alguno no lo conoces, escribe lo que te suena por el nombre. Es tu intuición; no se vale ver la siguiente actividad.
\end{softbox}

\checkbox Copié los 12 periféricos en mi cuaderno numerados.\\
\checkbox Al lado de cada uno, escribí para qué sirve con mis palabras.\\
\checkbox No miré la siguiente actividad ni busqué en internet.

\begin{infoband}{milcVerde}
\textbf{Lo que va al cuaderno (Actividad 1 · IDENTIFICA).} Título: «Actividad 1 --- Los 12 periféricos y mi intuición». \textbf{Formato:} lista numerada del 1 al 12. \textbf{Extensión:} 12 líneas. \textbf{Sabes que terminaste cuando} cada uno de los 12 periféricos tiene tu propia descripción al lado.
\end{infoband}

\puente{Después de la lista, aprendes la regla simple para clasificarlos.}

\begin{titlebox}{milcTurquesa}
Fase 2 · Sistematización con pensamiento computacional
\end{titlebox}

\begin{softbox}{milcTurquesa}{milcGris}{Cuatro pilares aplicados al tema}
Lo que acabaste de hacer ya es la base. Ahora vas a aprender \textbf{la regla simple} para clasificar cualquier periférico --- no solo los 12 de la lista, sino \emph{cualquiera que veas en el futuro}.
\end{softbox}

\small
\begin{tabularx}{\textwidth}{>{\bfseries\RaggedRight\arraybackslash}p{3.6cm}Y}
\rowcolor{milcTurquesa!90}\color{white}Pilar & \color{white}\bfseries Aplicación al tema\\
1. Descomposición & \textbf{Regla de oro --- 1 pregunta para clasificar.} Ante cualquier periférico, pregúntate: \emph{``¿Yo le doy algo al computador, o el computador me da algo a mí?''}. Si TÚ le das (tocas, escribes, hablas, mueves), es \textbf{entrada}. Si el computador te da (te muestra, te suena, te imprime, te vibra), es \textbf{salida}. Si pasan las 2 cosas en la misma pieza, es \textbf{mixto}.\\
2. Reconocimiento de patrones & \textbf{Los 6 periféricos de entrada típicos.} (1) \textbf{Teclado}: escribes letras y números. (2) \textbf{Ratón (mouse)}: mueves el cursor y haces clic. (3) \textbf{Micrófono}: hablas, el computador captura el sonido. (4) \textbf{Cámara web}: el computador captura imagen tuya. (5) \textbf{Escáner}: pone documentos en papel dentro del computador como imagen. (6) \textbf{Joystick}: en juegos, controla el personaje. \emph{Cómo se conectan:} casi todos por USB, algunos por Bluetooth (sin cable).\\
3. Abstracción & \textbf{Los 6 periféricos de salida típicos.} (1) \textbf{Monitor (pantalla)}: muestra todo lo que ves. (2) \textbf{Impresora}: pone lo digital en papel. (3) \textbf{Altavoz}: produce sonido al exterior. (4) \textbf{Audífonos}: igual que altavoz pero solo para ti. (5) \textbf{Proyector}: monitor gigante para muchas personas. (6) \textbf{Vibrador del celular}: te avisa con vibración. \emph{Cómo se conectan:} monitor por HDMI, audífonos por jack 3.5mm o Bluetooth, impresora por USB o wifi.\\
4. Algoritmo conceptual & \textbf{Los periféricos mixtos (entrada Y salida).} (1) \textbf{Pantalla táctil}: muestra (salida) y recibe tu toque (entrada). (2) \textbf{Audífonos con micrófono}: escuchas (salida) y hablas (entrada). (3) \textbf{Memoria USB}: lees archivos del computador (salida) y le pasas archivos (entrada). (4) \textbf{Módem o router wifi}: recibe internet (entrada) y manda a tus equipos (salida). \emph{Pista para reconocer mixtos:} si la pieza hace 2 cosas al mismo tiempo, probablemente es mixto.\\
\end{tabularx}
\normalsize

\begin{drawbox}{milcTurquesa}{La regla del 1 con 3 tipos}
\textbf{Regla:} ¿yo le doy o el computador me da? \\[1mm]
\textbf{Entrada:} teclado, ratón, micrófono, cámara, escáner, joystick. \\[1mm]
\textbf{Salida:} monitor, impresora, altavoz, audífonos, proyector, vibrador. \\[1mm]
\textbf{Mixto:} pantalla táctil, audífonos con micrófono, USB, módem.
\end{drawbox}

\begin{softbox}{milcMostaza}{milcGris}{Errores comunes y su corrección}
\textbf{Confusiones que se repiten cuando clasificas periféricos.} (1) \emph{``Los audífonos son siempre salida''} --- Falso. Los normales son salida, pero los \emph{audífonos con micrófono} (los que usas para juegos o llamadas) son mixtos. (2) \emph{``Una pantalla normal es entrada''} --- Falso. Una pantalla \emph{normal} es solo salida. Solo si es \emph{táctil} se vuelve mixta. (3) \emph{``La USB es periférico de almacenamiento''} --- Cierto pero incompleto. También es periférico mixto porque entra Y sale información. (4) \emph{``Sin cable significa que no es periférico''} --- Falso. Un periférico Bluetooth o wifi sigue siendo periférico, solo cambia la forma de conectarse.
\end{softbox}

\begin{infoband}{milcTurquesa}
\textbf{Lo que va al cuaderno (Actividad 2 · ANALIZA).} Título: «Actividad 2 --- Clasifico los 12 periféricos». \textbf{Formato:} ahora regresa a tu lista de la Actividad 1 y al lado de cada periférico, agrega \textbf{la etiqueta} ENTRADA, SALIDA o MIXTO. Si la intuición que escribiste antes coincide con la clasificación, ponte una palomita. Si no coincide, no la borres; déjala. \textbf{Extensión:} las mismas 12 líneas, ahora con la etiqueta. \textbf{Sabes que terminaste cuando} todos tienen su etiqueta.
\end{infoband}

\puente{Con la regla clara, construyes tu tabla maestra + diagrama con flechas.}

\begin{titlebox}{milcMagenta}
Fase 3 · Praxis con pensamiento computacional
\end{titlebox}

\begin{softbox}{milcMagenta}{milcGris}{Construyendo el producto con los cuatro pilares}
\textbf{Actividad 3 · CREA --- Tu tabla maestra + diagrama} (28 min · individual). Vas a hacer 2 cosas en una hoja completa del cuaderno. \textbf{Arriba (la mitad de la hoja):} la tabla maestra de los 12 periféricos. \textbf{Abajo (la otra mitad):} el diagrama con un computador en el centro y 6 periféricos conectados con flechas.
\end{softbox}

\small
\begin{tabularx}{\textwidth}{>{\bfseries\RaggedRight\arraybackslash}p{3.6cm}Y}
\rowcolor{milcMagenta!90}\color{white}Pilar & \color{white}\bfseries Aplicación al producto\\
Descomposición & \textbf{Paso 1 --- Tabla maestra (mitad superior).} Dibuja una tabla de \textbf{12 filas y 4 columnas}. Encabezados: \emph{Nombre}, \emph{Tipo} (entrada / salida / mixto), \emph{Cómo se conecta} (USB, HDMI, jack, Bluetooth, wifi), \emph{Ejemplo de uso real}. Llena las 12 filas con los periféricos de la Actividad 1. En la columna ``ejemplo'', pon una situación tuya: \emph{``Lo uso para...''}, no copies del libro.\\
Patrones & \textbf{Paso 2 --- Diagrama del computador con periféricos (mitad inferior).} Dibuja en el centro de la mitad inferior un \textbf{rectángulo grande} y adentro escribe \textbf{``COMPUTADOR''}. Ese es el centro de tu diagrama.\\
Abstracción & \textbf{Paso 3 --- Periféricos de entrada (3 a la izquierda).} A la izquierda del rectángulo dibuja 3 íconos pequeños y etiquétalos: \textbf{teclado}, \textbf{ratón}, \textbf{micrófono}. Desde cada uno traza una \textbf{flecha apuntando hacia el rectángulo} (porque la información va DEL periférico AL computador).\\
Algoritmo de armado & \textbf{Paso 4 --- Periféricos de salida (3 a la derecha).} A la derecha del rectángulo dibuja otros 3 íconos: \textbf{monitor}, \textbf{altavoz}, \textbf{impresora}. Desde el rectángulo traza \textbf{flechas saliendo hacia ellos} (porque la información va DEL computador AL periférico).\\
\end{tabularx}
\normalsize

\begin{drawbox}{milcMagenta}{Antes de mostrar tu hoja}
\checkbox La tabla maestra tiene 12 filas y 4 columnas. \\[1mm]
\checkbox Cada periférico tiene su tipo, conexión y ejemplo de uso. \\[1mm]
\checkbox El diagrama tiene un computador en el centro. \\[1mm]
\checkbox Hay 3 periféricos de entrada a la izquierda con flechas hacia el centro. \\[1mm]
\checkbox Hay 3 periféricos de salida a la derecha con flechas hacia afuera. \\[1mm]
\checkbox Está firmado con nombre y fecha.
\end{drawbox}

\begin{softbox}{milcMagenta}{milcGris}{Plantilla de trabajo}
\textbf{Mitad superior:} tabla 12 × 4 (nombre, tipo, conexión, ejemplo). \textbf{Mitad inferior:} rectángulo central con ``COMPUTADOR''. Izquierda: 3 periféricos de entrada con flechas hacia el rectángulo, etiqueta ENTRADA arriba. Derecha: 3 periféricos de salida con flechas saliendo, etiqueta SALIDA arriba.
\end{softbox}

\begin{infoband}{milcMagenta}
\textbf{Lo que va al cuaderno (Actividad 3 · CREA).} Título: «Actividad 3 --- Mi tabla maestra de periféricos + diagrama». \textbf{Formato:} tabla 12 × 4 arriba + diagrama con flechas abajo. \textbf{Extensión:} 1 hoja entera. \textbf{Sabes que terminaste cuando} podrías mostrarle la hoja a un primo y él entendería en 3 minutos qué es entrada, salida y mixto.
\end{infoband}

\puente{El producto es la tabla maestra firmada + el diagrama con 6 periféricos conectados.}

\begin{titlebox}{milcMostaza}
Producto de la sesión
\end{titlebox}
\textbf{Tabla maestra de 12 periféricos + diagrama con 6 flechas · firmado}.
\begin{softbox}{milcMostaza}{milcGris}{Criterios de calidad}
(1) La tabla cubre los 12 periféricos con sus 4 datos. (2) La columna 'tipo' tiene clasificación correcta (entrada/salida/mixto). (3) El diagrama tiene 3 periféricos de entrada con flechas hacia el computador. (4) El diagrama tiene 3 periféricos de salida con flechas hacia afuera. (5) Está firmado y fechado.
\end{softbox}

\puente{Antes de salir, verificas que cada periférico tenga sus 4 datos completos en la tabla.}

\begin{titlebox}{milcVino}
Fase 4 · Evaluación liberadora — cinco dimensiones
\end{titlebox}

\begin{softbox}{milcVino}{milcGris}{Sentido de la evaluación}
Evaluar no es solo revisar una nota. Es reconocer qué aprendí, qué cambió en mí, qué emociones manejé, cómo me relaciono con la ciudad y la comunidad, y qué le dejaré a quien viene detrás.
\end{softbox}

\textbf{Mi reflexión liberadora en cinco dimensiones.} Completa cada frase iniciada:

\small
\begin{tabularx}{\textwidth}{>{\bfseries\RaggedRight\arraybackslash}p{3.4cm}Y>{\RaggedRight\arraybackslash}p{4.6cm}}
\rowcolor{milcVino}\color{white}Dimensión & \color{white}\bfseries Mi respuesta (frame starter) & \color{white}\bfseries Algo que descubrí\\
1. Personal & Lo que más cambió en mí fue \linead{6.5cm} & \linead{4.2cm}\\
2. Emocional & La emoción más fuerte fue \linead{2.5cm} y la regulé \linead{3cm} & \linead{4.2cm}\\
3. Ciudadana & En democracia esto importa porque \linead{6cm} & \linead{4.2cm}\\
4. Local (Cartago) & En mi barrio sería útil para \linead{6.5cm} & \linead{4.2cm}\\
5. Intergeneracional & A mi abuela/abuelo le diría \linead{6.5cm} & \linead{4.2cm}\\
\end{tabularx}
\normalsize

\begin{infoband}{milcVino}
\textbf{Para la próxima sustentación.} Vuelve a esta tabla en seis meses. Verás que las respuestas cambian — no porque lo aprendido cambie, sino porque tú cambias. La evaluación liberadora es una conversación contigo a través del tiempo.
\end{infoband}


\puente{Cerramos con 3 voces que te ayudan a entender por qué conocer las herramientas te hace mejor usuario.}

\begin{titlebox}{milcGradeDark}
Triángulo de pensamiento · cierre filosófico
\end{titlebox}

\begin{softbox}{milcVino}{milcGris}{Dussel · lente del nosotros}
\textit{````Conocer las herramientas a la mano es conocer cómo vivimos.''''}

Cada periférico es una \textbf{decisión cultural}: alguien decidió que el teclado tuviera esas teclas, que el ratón fuera ese tamaño, que la pantalla mostrara colores. \textbf{Conocer los periféricos es entender el mundo que nosotros construimos}. No son neutros: son fruto de decisiones humanas, y por eso pueden cambiar (mira los celulares actuales --- ya casi no se usan teclas físicas).

\textbf{¿Cuál es el periférico que uso más cada día? ¿Lo escogí yo o me lo asignaron?}
\end{softbox}
\linead{\linewidth}\\[1mm]
\linead{\linewidth}

\begin{softbox}{milcMostaza}{milcGris}{Estoicismo · Marco Aurelio · cuidado interior}
\textit{````Cuidar la herramienta es prolongar su utilidad y la nuestra.''''}

Los periféricos son \textbf{los que más se dañan} en un computador: el teclado por golpes, el ratón por suciedad, la impresora por mal uso, los audífonos por estirones. Si los \textbf{cuidas} (limpieza, posición, no halarlos por el cable), te duran años. Si los maltratas, los gastas en meses. \emph{Cuidar las herramientas es respetar el trabajo que harán contigo}.

\textbf{¿Trato mis periféricos como herramientas valiosas o como cosas desechables que pronto reemplazaré?}
\end{softbox}
\linead{\linewidth}\\[1mm]
\linead{\linewidth}

\begin{softbox}{milcTurquesa}{milcGris}{Floridi · lente de la infoesfera}
\textit{````El cuerpo de cada usuario se extiende a través de sus periféricos.''''}

Cuando usas el teclado, tus dedos se \textbf{extienden} hacia los caracteres en la pantalla. Cuando usas un micrófono, tu voz se \textbf{extiende} al lugar al que llamas. Cuando usas Skype con cámara, tu \emph{presencia} se extiende hasta otra ciudad. Los periféricos \textbf{prolongan tu cuerpo}. Por eso, en cierto sentido, son parte de ti cuando los usas. Cuidarlos es cuidarte.

\textbf{¿De qué manera mis periféricos me hacen \emph{más yo}, y de qué manera me alejan de mí mismo?}
\end{softbox}
\linead{\linewidth}\\[1mm]
\linead{\linewidth}

\begin{titlebox}{milcVino}
Compromiso final
\end{titlebox}

\small
\begin{tabularx}{\textwidth}{>{\RaggedRight\arraybackslash}p{3.0cm}YYY}
\rowcolor{milcVino}\color{white}\bfseries Criterio & \color{white}\bfseries Lo logré & \color{white}\bfseries En proceso & \color{white}\bfseries Necesito apoyo\\
Comprensión del tema & Explico ideas centrales con mis palabras. & Reconozco ideas pero falta precisión. & Necesito repasar conceptos base.\\
Pensamiento computacional & Apliqué los 4 pilares al diseñar mi producto. & Apliqué algunos pilares. & Necesito releer la fase 2.\\
Aplicación y producto & Mi producto cumple los criterios y se puede revisar. & Producto funcional con ajustes pendientes. & Aún en construcción.\\
Reflexión 5 dimensiones & Respondí cada dimensión con honestidad. & Respondí algunas con detalle. & Mis respuestas son muy generales.\\
Triángulo de pensamiento & Conecté las 3 voces con mi trabajo real. & Conecté al menos una. & No relaciono las voces con mi proceso.\\
\end{tabularx}
\normalsize

\textbf{Hoy entendí que:}\\[1mm]
\linead{\linewidth}\\[1mm]
\linead{\linewidth}

\textbf{Mi compromiso para la próxima sesión es:}\\[1mm]
\linead{\linewidth}\\[1mm]
\linead{\linewidth}

\begin{softbox}{milcMostaza}{milcGris}{Cita de despedida · Marco Aurelio}
\textit{``El obstáculo en el camino se vuelve el camino. Lo que estorba al camino se vuelve camino.''}\\[1mm]
Cada error de hoy, bien observado, se convierte en pieza del camino que pisarás mañana.
\end{softbox}

\small
\begin{tabularx}{\textwidth}{>{\bfseries}p{3.6cm}Y}
\rowcolor{milcGradeDark!12}Nombre completo: & \linead{10cm}\\
Fecha: & \linead{4cm} \quad Curso/grupo: \linead{4cm}\\
Mi firma: & \linead{9cm}\\
\end{tabularx}
\normalsize

\begin{infoband}{milcGradeDark}
\textbf{Para la próxima sesión.} Esta semana, observa cuántos periféricos usas cada día y por qué tipo de conexión. ¿Tu teclado es USB o Bluetooth? ¿Tu impresora va por wifi o cable? La próxima clase aprendes la \textbf{diferencia entre hardware y software}: el cuerpo y las instrucciones del computador. Hoy ya conoces el cuerpo (gabinete + periféricos); falta entender las instrucciones que lo mueven.
\end{infoband}

\end{document}
